\documentclass[11pt]{article}
\usepackage{setspace}    
\onehalfspacing
\usepackage{graphicx}    
\usepackage{amsmath}     
%\usepackage{natbib}      % Apagado en esta tarea
\usepackage{array}       
\usepackage{multirow}    
\usepackage{siunitx}     
\usepackage{textcomp}    
\usepackage[utf8]{inputenc}  %Apagado, motor incluye utf8
\usepackage[spanish,es-tabla]{babel}  
\usepackage{csquotes}   
\usepackage[top=1in,right=1in,left=1in,bottom=1in]{geometry} 
\usepackage{makecell}
\usepackage{hyperref}        

\hypersetup{
    unicode=true,            
    pdftoolbar=true,         
    pdfmenubar=true,         
    pdffitwindow=false,      
    pdfstartview={FitH},     
    pdftitle={},             
    pdfauthor={},            
    pdfsubject={},           
    pdfcreator={},           
    pdfproducer={},          
    pdfkeywords={},          
    pdfnewwindow=true,       
    colorlinks=true,         
    linkcolor=blue,          
    citecolor=blue,          
    filecolor=blue,          
    urlcolor=blue            
}

\usepackage[backend=biber,style=chicago-authordate,natbib=true,language=spanish]{biblatex} %paquete para manejar citas
\addbibresource{Consti.bib} %Documento que contiene anotaciones bibliográficas

\title{SOCI 4186-00N\\ Tarea \textnumero 2 - Ejercicio de educación cívica}

\author{Pon tu nombre aquí}
\date{\today}

\begin{document}
\singlespacing
\maketitle
\onehalfspacing
Fecha límite de entrega: 14 de febrero de 2025.

\begin{enumerate}
    \item Descargarán este documento en el GitHub de la clase, y lo subirán a Overleaf.
    \item Cambiarán el nombre que sale arriba (al final del preámbulo, poco antes del \texttt{\textbackslash begin\{document\}}, donde dice \texttt{\textbackslash author\{Pon tu nombre aquí\}} y pondrán en su lugar su nombre(s) y apellidos.
    \item Prestarán atención a la ortografía de la lengua española, asegurándose de colocar acentos, diéresis, o tildes donde deban ir. Para esto, tienen dos opciones 
    \begin{enumerate}
        \item si su teclado estuviera en español, usen los símbolos donde es menester,
        \item si su teclado no estuviera en español, está la opción indicada en clase de usar \texttt{\textbackslash} y el símbolo necesario (por ejemplo, \textbackslash' o \textbackslash'{} para acento (\'a o \'{a}), \textbackslash\"\ o \textbackslash\"{} para diéresis (\"u o \"{u}), \textbackslash\textasciitilde\ para \~n)
        \item Alternativamente, si prefiriera escribir en inglés \underline{la tarea entera}, podrá hacerlo, ciñéndose estrictamente a una de las tres ortografías principales de ese idioma, es decir, inglés británico, inglés Oxford, o inglés estadounidense.
    \end{enumerate}
    \item Completarán como mejor sea posible la tarea, y podrán, si entienden necesario hacerlo, trabajar la tarea con compañeres de clase. Si así lo hicieren, pondrán antes de la sección 1"A una oración indicándome con quién(es) trabajó la tarea. De lo contrario, dejarán la línea que está actualmente escrita.
    \item  Una vez terminen, apretarán el símbolo de descarga (a la derecha del botón de Compilar, con flecha descendiente), y descargarán el PDF para enviarlo a mi persona.
    \item Someterán la tarea mediante la aplicación de Teams el viernes 14 de febrero.
\end{enumerate}
\textbf{Declaro que soy la única persona que trabajó en solucionar esta tarea.}

\section{Viviendo en democracia}
A menudo escuchamos sobre términos políticos de uso común, pero no abundamos en profundidad sobre los mismos, y carecemos por ello de comprensión de los sistemas sociales complejos en los que navegamos. En esta tarea, aplicarán el conocimiento adquirido de procesamiento de textos en \LaTeX en las pasadas semanas para profundizar en su conocimiento de sus derechos, así como hablar sobre eventos actuales de relevancia para ustedes.

\textbf{Pregunta 1"A}
Seleccione uno de los derechos plasmados en el Artículo II \citep{ConstitucionPR} que le resulte particularmente relevante para la convivencia social (por ejemplo, libertad de expresión, derecho a la educación, derecho a la huelga, etc.). 

\begin{enumerate}
    \item Explique brevemente, en sus propias palabras, en qué consiste ese derecho y cómo la Constitución de Puerto Rico lo establece.
    \item Cite al menos \textbf{un} artículo académico, libro o fuente no académica (por ejemplo, un artículo periodístico) que discuta la evolución o los retos actuales de ese derecho. Debe usar bibliografía en formato discutido en clase.
    \item Reflexione brevemente sobre la importancia de este derecho en la vida diaria, conectándolo con algún ejemplo reciente o de su experiencia personal.
\end{enumerate}

\textbf{Pregunta 2"A}

La Asamblea Multisectorial del Recinto de Río Piedras ocurre esta semana, el martes 11 de febrero de 2025. Citando dos medios de prensa distintos (periódico --impreso o digital--, radio o televisión) indique qué se determinó en la misma. Use cita parentética o textual, dependiendo del contexto adecuado.


\printbibliography

\end{document}